
\subsection{Struct}

\lstinputlisting{geo/base.cpp}

\iffalse
\subsection{Convex Hull}
\subsubsection{Graham Scan}

\begin{lstlisting}[language=C++] 
using Point= complex<int>;

inline long long crossProduct(const Point& vec1, const Point& vec2) {
    return (long long)vec1.real() * vec2.imag() - vec1.imag() * vec2.real();
}

int orientation(Point p, Point q, Point r) {

    Point vec1= p-q;

    Point vec2= q-r;

    int val = crossProduct(vec1,vec2);

    if (val == 0) return 0;  // Collinear
    return (val > 0) ? 1 : 2; // Clockwise or counterclockwise
}

bool comparePoints(const Point& p0, const Point& p1, const Point& p2) {
    int orient = orientation(p0, p1, p2);

    if (orient == 0)
        return (norm(p1 - p0) < norm(p2 - p0)); // Sort by distance from p0

    return (orient == 2); // Sort counterclockwise
}

vector<Point> convexHull(vector<Point>& points){
    if(points.size()<3){
        return points;
    }

    auto pivot = *min_element(points.begin(), points.end(), [](const Point& a, const Point& b) {
        return (imag(a) == imag(b)) ? (real(a) < real(b)) : (imag(a) < imag(b)); });
    
    swap(points[0],pivot);

    sort(points.begin()+1,points.end(),[&points](const Point& a, const Point& b){
        return comparePoints(points[0],a,b);
    } );



    vector<Point> hull({points[0],points[1],points[2]});
    
    for(size_t i=3; i<points.size(); i++){
        while(hull.size()>2 and orientation(hull[hull.size()-2],hull.back(),points[i])!=2){
            hull.pop_back();
        }
        hull.push_back(points[i]);   
    }

    return hull;
}

\end{lstlisting}
\fi